\subsection{An explicit approximation class and a uniform scale}

\begin{definition}[Coherent pair approximation on a partition]
    \label{def:ldp_coherent_pair_approximable}
    \lean{MPSPreparation.IsCoherentPairApproximable}
    \leanok
    \uses{def:ldp_chain_pos_state}
    A nonzero inhomogeneous chain state $\phi_N$ is approximable on a given
    partition by coherent pair products with $b$ labels and error $\delta$ if
    there are pairs and amplitudes as in
    (\ref{eq:ldp_inhomogeneous_long_range_pairs}) such that
    $1-|\langle\Omega',\phi_{\rm pos}/\|\phi_{\rm pos}\|\rangle|\le\delta$.
    This specifies a class of approximation hypotheses for the remark
    after (19) in~\cite{Malz2024Preparation}; that remark does not give a
    separate definition of the inhomogeneous long-range class.
\end{definition}

\begin{lemma}[Physical preparation under coherent pair approximation]
    \label{lem:ldp_coherent_pair_approximation_preparation}
    \lean{MPSPreparation.exists_isPreparedWithMeasurementsAndCircuitInDepth_of_isCoherentPairApproximable}
    \leanok
    \uses{def:ldp_coherent_pair_approximable,
        lem:ldp_inhomogeneous_long_range_preparation}
    Fix $d\ge2$, $D\ge1$ and $b$. There are $C,L_0$ such that a coherent
    pair approximation with error $\delta$ on blocks of lengths
    $L_0\le\ell_k\le L$ gives a unit physical state prepared with measurements
    followed by a circuit in depth $CL$ and error at most $\delta$ against
    the normalized chain state.
\end{lemma}
\begin{proof}\leanok
    Apply Lemma~\ref{lem:ldp_inhomogeneous_long_range_preparation} to the pairs
    and amplitudes in the approximation. Its exact overlap identity transfers
    the assumed error unchanged.
\end{proof}

\begin{lemma}[Uniform exponential coherent-pair rate gives a logarithmic scale]
    \label{lem:ldp_coherent_pair_uniform_rate}
    \lean{MPSPreparation.exists_inhomogeneousLongRange_depth_le_log_of_uniform_rate}
    \leanok
    \uses{lem:ldp_coherent_pair_approximation_preparation}
    Fix $d\ge2$, $D\ge1$, $b$, $K,r>0$ and an integer $k\ge0$. Let $A_N$
    be a family of inhomogeneous chains with bond dimension $D$. Suppose
    that for every $N\ge1$ and $1\le q\le N$ there is a partition into blocks
    with $q\le\ell_j\le2q$ on which the chain has a coherent pair
    approximation with $b$ labels and error at most $KN^k e^{-rq}$.
    There is $C\ge0$, independent of $N$ and $\varepsilon$, such that for
    every $N\ge1$ and $0<\varepsilon\le1$ a unit state of error at most
    $\varepsilon$ is prepared with measurements followed by a circuit in
    depth at most $C(1+\log(N/\varepsilon))$.
\end{lemma}
\begin{proof}\leanok
    Let $C_b,L_0$ be the constants of
    Lemma~\ref{lem:ldp_coherent_pair_approximation_preparation}. Set
    $a=\max(k,1)/r$, $B=\max(\log K,0)/r+L_0+1$ and
    $q=\lceil a\log(N/\varepsilon)+B\rceil$. Then $q\ge L_0$,
    $KN^ke^{-rq}\le\varepsilon$ and
    $q\le(a+B+1)(1+\log(N/\varepsilon))$. For $q\le N$, apply the
    approximation hypothesis and the preceding lemma. For $q>N$, exact
    preparation of the entire inhomogeneous chain has depth at most $C_eN$,
    which is bounded by $C_eq$. Thus
    $C=(2C_b+C_e)(a+B+1)$ suffices in both cases. The additive one includes
    a one-site ring at accuracy one.
\end{proof}

\begin{lemma}[The label count is bounded by the pair dimension]
    \label{lem:ldp_coherent_pair_label_bound}
    \lean{MPSPreparation.IsCoherentPairApproximable.label_le}
    \leanok
    \uses{def:ldp_coherent_pair_approximable}
    On a nonempty partition, every coherent pair approximation with
    bond dimension $D$ has at most $D^2$ labels.
\end{lemma}
\begin{proof}\leanok
    At one window the label vectors are orthonormal in a space of
    dimension $D^2$. The corresponding matrix is an isometry, so its
    number of columns does not exceed its number of rows.
\end{proof}

\begin{lemma}[Uniform constants with varying label counts]
    \label{lem:ldp_coherent_pair_varying_labels}
    \lean{MPSPreparation.exists_coherentPairPreparation_uniform_labels,
        MPSPreparation.exists_inhomogeneousLongRange_depth_le_log_of_uniform_rate_varying_labels}
    \leanok
    \uses{lem:ldp_coherent_pair_label_bound,
        lem:ldp_coherent_pair_approximation_preparation}
    For fixed $d\ge2$ and $D\ge1$, the constants $C,L_0$ for coherent
    pair preparation can be chosen independently of the number of
    labels. In the uniform exponential-rate hypothesis, the number
    of labels may depend on the length and the chosen partition.
    The resulting depth bound remains $C'(1+\log(N/\varepsilon))$,
    with constants independent of all these label counts.
\end{lemma}
\begin{proof}\leanok
    Choose preparation constants for each of the finitely many label
    counts from $0$ through $D^2$. Their sums bound all choices at once.
    Use these uniform constants in the blocking argument of
    Lemma~\ref{lem:ldp_coherent_pair_uniform_rate}.
\end{proof}

\subsection{What qualitative convergence does not imply}

\begin{definition}[Delayed blocking error]
    \label{def:ldp_delayed_blocking_error}
    \lean{MPSPreparation.delayedBlockingError}
    \leanok
    For integers $N,q\ge0$, define $\delta(N,q)=1/(N+1)$ if $q<N$ and
    $\delta(N,q)=0$ if $q\ge N$.
\end{definition}

\begin{lemma}[Qualitative convergence with exact whole-ring blocking]
    \label{lem:ldp_delayed_blocking_convergence}
    \lean{MPSPreparation.delayedBlockingError_bounds,
        MPSPreparation.antitone_delayedBlockingError,
        MPSPreparation.delayedBlockingError_eq_zero,
        MPSPreparation.tendsto_delayedBlockingError}
    \leanok
    \uses{def:ldp_delayed_blocking_error}
    The error $\delta(N,q)$ is nonincreasing in $q$, belongs to
    $[0,1/(N+1)]$ and vanishes for $q\ge N$. For every schedule $q_N$,
    including any logarithmic schedule, $\delta(N,q_N)\to0$ as $N\to\infty$.
\end{lemma}
\begin{proof}\leanok
    The assertions follow from the two cases in the definition and the
    bound $0\le\delta(N,q_N)\le1/(N+1)$.
\end{proof}

\begin{lemma}[No uniform logarithmic accuracy scale from these errors]
    \label{lem:ldp_delayed_blocking_no_uniform_scale}
    \lean{MPSPreparation.exists_delayedBlockingError_gt_of_log_bound,
        MPSPreparation.not_exists_uniform_log_blocking_for_delayedBlockingError}
    \leanok
    \uses{lem:ldp_delayed_blocking_convergence}
    For every $C\ge0$ there are arbitrarily large $N\ge2$ and an accuracy
    $0<\varepsilon\le1$ such that every integer
    $q\le C\log(N/\varepsilon)$ satisfies $\delta(N,q)>\varepsilon$.
    Consequently there is no uniform logarithmic accuracy-dependent blocking
    coefficient for these error functions.
\end{lemma}
\begin{proof}\leanok
    Choose $\varepsilon=1/(2(N+1))$. Since
    $\log(N/\varepsilon)\le3\log(N+1)$ and
    $3C\log(N+1)\le(N+1)/2<N$ for all sufficiently large $N$,
    every permitted $q$ is less than $N$. Its error is then
    $1/(N+1)>\varepsilon$.
\end{proof}

This example concerns numerical error information alone. It does not
construct matrix product states whose optimal errors are $\delta(N,q)$,
and it gives no circuit-depth lower bound. A theorem deriving a uniform
scale for a class of matrix product states must use additional information
about that class.

\subsection{Coordinate support in the tree construction}

\begin{definition}[Rotated GHZ flags]
    \label{def:ldp_rotated_ghz_flags}
    \lean{MPSTensor.flagRotation, MPSTensor.rotatedGHZFlag}
    \leanok
    Let $R=\left(\begin{smallmatrix}3/5&-4/5\\4/5&3/5\end{smallmatrix}\right)$
    and $B^j=R\ket j\bra j$ for $j=0,1$.
\end{definition}

\begin{lemma}[The rotation and its physical range are isometric]
    \label{lem:ldp_rotated_ghz_coisometry}
    \lean{MPSTensor.conjTranspose_flagRotation_mul,
        MPSTensor.physicalMatrix_rotatedGHZFlag_mul_conjTranspose}
    \leanok
    \uses{def:ldp_rotated_ghz_flags}
    The matrix $R$ satisfies $R^\dagger R=I$. The physical matrix
    $F_{j,(a,b)}=(B^j)_{ab}$ satisfies $FF^\dagger=I_2$.
\end{lemma}
\begin{proof}\leanok
    Multiply the displayed rational matrices.
\end{proof}

\begin{lemma}[The rotated flags have no injective coordinate support]
    \label{lem:ldp_rotated_ghz_no_coordinate_support}
    \lean{MPSTensor.not_exists_isInjectiveOn_rotatedGHZFlag}
    \leanok
    \uses{lem:ldp_rotated_ghz_coisometry}
    There is no set $S$ of virtual coordinate pairs such that the matrices
    $B^j$ vanish outside $S$ and span all matrix units indexed by $S$.
\end{lemma}
\begin{proof}\leanok
    Since $(B^0)_{00}=3/5\ne0$, such a set would contain $(0,0)$, so its
    matrix span would contain $E_{00}$. Every matrix in that span satisfies
    $3X_{10}-4X_{00}=0$, whereas $E_{00}$ does not.
\end{proof}

Thus bounded rank and an isometric physical range do not imply the
coordinate-support hypothesis of
Lemma~\ref{thm:ldp_coherent_supported_blocks}. A general inhomogeneous
tree argument requires virtual subspaces and vertex maps depending on their
positions, or a proved gauge reduction to the supported setting. This
obstruction to that hypothesis is not an impossibility result for such
general trees.

\begin{lemma}[Full-domain isometric polar completion]
    \label{lem:ldp_isometric_polar_completion}
    \lean{Matrix.exists_isometry_mul_polarPos_eq}
    \leanok
    Let $B\colon\C^p\to\C^m$ with $p\le m$, and let
    $P=(B^\dagger B)^{1/2}$. There is an isometry $X$ with $XP=B$,
    without an injectivity assumption on $B$.
\end{lemma}
\begin{proof}\leanok
    Choose an isometric inclusion $E\colon\C^p\to\C^m$.
    The matrices $B$ and $EP$ have the same Gram matrix $P^2$, so a unitary
    $U$ on the output space satisfies $B=UEP$. Then $X=UE$ suffices.
\end{proof}

\begin{lemma}[Isometric completion inside the children of a binary node]
    \label{lem:ldp_isometric_polar_merge}
    \lean{Matrix.exists_isometry_polarMerge_of_range,
        MPSPreparation.exists_isometry_binaryPolarMerge}
    \leanok
    \uses{lem:ldp_isometric_polar_completion}
    Let $J\colon\C^n\to\C^m$ be an isometry, let
    $B\colon\C^p\to\C^m$ satisfy $JJ^\dagger B=B$, and assume $p\le n$.
    There is an isometry $W\colon\C^p\to\C^n$ such that $JW$ is an
    isometry and $(JW)(B^\dagger B)^{1/2}=B$.
    In particular, for two child isometries
    $L\colon\C^\chi\to\mathcal H_1$ and
    $R\colon\C^\chi\to\mathcal H_2$ with $\chi\ge1$, this applies to
    $J=L\otimes R$ and every compressed parent matrix
    $B=JC$, $C\colon\C^\chi\to\C^{\chi^2}$.
\end{lemma}
\begin{proof}\leanok
    Put $C=J^\dagger B$. The range assumption gives
    $C^\dagger C=B^\dagger B$. Complete the polar factor of $C$ by
    Lemma~\ref{lem:ldp_isometric_polar_completion}; the equality of the
    positive factors and $JC=B$ give the claim. The tensor product of the
    child maps is an isometry and $\chi\le\chi^2$.
\end{proof}
